\documentclass[9pt,twocolumn]{article}

\usepackage[margin=0.58in,columnsep=0.22in]{geometry}
\usepackage{amsmath,amssymb,mathtools}
\usepackage{microtype}
\usepackage{newtxtext,newtxmath}

\setlength{\parindent}{0.8em}
\setlength{\parskip}{0pt}
\setlength{\abovedisplayskip}{4pt}
\setlength{\belowdisplayskip}{4pt}
\setlength{\abovedisplayshortskip}{2pt}
\setlength{\belowdisplayshortskip}{2pt}

\title{\vspace{-2.4em}\bfseries
\texttt{sin()} Stays \texttt{sin()}; Angles Become Plural
\vspace{-0.4em}}

\author{\small Wilder Blair Munro $\times$ Aurora}
\date{}

\begin{document}
\maketitle
\vspace{-2.2em}

\begin{abstract}
\noindent
For $\beta=v/c$, special relativity admits simultaneously

$$
\beta=\sin\theta=\tanh\chi,
\qquad
\gamma=\sec\theta=\cosh\chi,
$$

where $\theta=\operatorname{gd}\chi$ is the Gudermannian relation between circular angle and rapidity. Define additionally

$$
\alpha=\beta^{-1},\qquad
\delta=(1-\beta^2)^{-1/2}.
$$

Then

$$
\alpha=\csc\theta=\coth\chi,\qquad
\delta=\sec\theta=\cosh\chi.
$$

Nothing has happened to \texttt{sin()}. What changed was the representation of angle. We propose treating circular tilt, hyperbolic rapidity, and reciprocal aperture as coordinate bodies of a single \emph{vantage-angle object}. Under this interpretation a lift--rotate--squeeze construction is not motivated by preference for Euclidean rotations; it is the explicit realization of a change of angular chart.
\end{abstract}

\vspace{-0.6em}
\section*{I. One motion, several angles}

Let

$$
\beta_v=\frac{v}{c},\qquad
\beta_c=\sqrt{1-\beta_v^2},
$$

so that
\begin{equation}
\beta_v^2+\beta_c^2=1.
\tag{1}
\end{equation}
There therefore exists $\theta\in(-\pi/2,\pi/2)$ satisfying
\begin{equation}
\beta_v=\sin\theta,\qquad
\beta_c=\cos\theta.
\tag{2}
\end{equation}

Special relativity instead naturally parameterizes the same speed by rapidity
\begin{equation}
\beta_v=\tanh\chi,
\qquad
\gamma_v=\cosh\chi,
\qquad
\gamma_v^{-1}=\operatorname{sech}\chi.
\tag{3}
\end{equation}
Equations (2)--(3) imply
\begin{equation}
\sin\theta=\tanh\chi,\qquad
\cos\theta=\operatorname{sech}\chi.
\tag{4}
\end{equation}
This correspondence is not novel. It is exactly the Gudermannian:
\begin{equation}
\boxed{\theta=\operatorname{gd}\chi.}
\tag{5}
\end{equation}

Thus the circular and hyperbolic descriptions require no complexification:

\begin{equation*}
\begin{array}{c|c}
\text{circular} & \text{hyperbolic}\\ \hline
\sin\theta & \tanh\chi\\
\cos\theta & \operatorname{sech}\chi\\
\tan\theta & \sinh\chi\\
\csc\theta & \coth\chi\\
\sec\theta & \cosh\chi
\end{array}
\tag{6}
\end{equation*}

The physical invariant is not either parameterization. It is the relation they coordinatize.

Define a \emph{vantage angle}

$$
\mathfrak a\in\mathcal A
$$

with circular and hyperbolic charts

$$
q_\theta(\mathfrak a)=\theta,
\qquad
q_\chi(\mathfrak a)=\chi,
$$

related by (5). Then
\begin{equation}
\beta_v
=\sin(q_\theta\mathfrak a)
=\tanh(q_\chi\mathfrak a).
\tag{7}
\end{equation}

The word ``angle'' has become plural without modifying either function.

\section*{II. Reciprocal angles}

Now define
\begin{equation}
\alpha:=\frac1{\beta_v}=\frac{c}{v},
\qquad
\delta:=\frac1{\beta_c}.
\tag{8}
\end{equation}
From (2),
\begin{equation}
\boxed{
\alpha=\csc\theta,\qquad
\delta=\sec\theta.
}
\tag{9}
\end{equation}
From (6),
\begin{equation}
\boxed{
\alpha=\coth\chi,\qquad
\delta=\cosh\chi.
}
\tag{10}
\end{equation}

Consequently

\begin{equation*}
\delta=\gamma_v.
\tag{11}
\end{equation*}

If a complementary Lorentz factor is defined by

$$
\gamma_c:=(1-\beta_c^2)^{-1/2},
$$

then (1) gives

$$
\gamma_c=\frac1{|\beta_v|}.
$$

For $\beta_v>0$,
\begin{equation}
\boxed{
\alpha=\gamma_c,\qquad
\delta=\gamma_v.
}
\tag{12}
\end{equation}

The reciprocal pair obeys
\begin{equation}
\boxed{
\frac1{\alpha^2}+\frac1{\delta^2}=1,
}
\tag{13}
\end{equation}
or equivalently
\begin{equation}
\boxed{
\alpha^2+\delta^2=\alpha^2\delta^2.
}
\tag{14}
\end{equation}
Hence
\begin{equation}
\sqrt{\alpha^2+\delta^2}=\alpha\delta
\tag{15}
\end{equation}
in the positive quadrant.

Equation (15) is the reciprocal triangle.

\section*{III. The triangle survives inversion}

Write the normalized velocity triangle as
\begin{equation}
T_\beta=(\beta_c,\beta_v,1).
\tag{16}
\end{equation}
Define
\begin{equation}
\kappa
:=
\frac1{\beta_v\beta_c}
=\alpha\delta.
\tag{17}
\end{equation}
Uniform scaling gives

$$
\kappa T_\beta
=
\left(
\frac1{\beta_v},
\frac1{\beta_c},
\frac1{\beta_v\beta_c}
\right),
$$

therefore
\begin{equation}
\boxed{
\kappa T_\beta=(\alpha,\delta,\alpha\delta)=:T_\rho.
}
\tag{18}
\end{equation}

The ``reciprocal triangle'' is consequently not a second arbitrary triangle. It is similar to the first.

Indeed,
\begin{equation}
\frac{\delta}{\alpha}
=\frac{\beta_v}{\beta_c}
=\tan\theta.
\tag{19}
\end{equation}
and
\begin{equation}
\frac{\alpha}{\delta}
=\frac{\beta_c}{\beta_v}
=\cot\theta.
\tag{20}
\end{equation}

The direction survives. Only scale changes.

Thus the sequence

$$
(\beta_c,\beta_v,1)
\longrightarrow
(\alpha,\delta,\alpha\delta)
$$

is exactly

\begin{equation}
\boxed{\text{same angle; reciprocal body.}}
\tag{21}
\end{equation}

This is the geometric fact hidden by the original notational overload.

\section*{IV. The apparent sin}

The expression
\begin{equation}
\sin\alpha=\frac1\alpha
\tag{22}
\end{equation}
is false if $\alpha\in\mathbb R$ is interpreted as an ordinary circular angle measured in radians.

But (9) says something more precise:

\begin{equation}
\alpha=\csc\theta
\quad\Longrightarrow\quad
\sin\theta=\frac1\alpha.
\tag{23}
\end{equation}

The original abuse therefore conflates an object with one of its coordinates.

Let the reciprocal chart be

$$
q_\rho(\mathfrak a)=(\alpha,\delta).
$$

The intrinsic circular projection

$$
S:\mathcal A\rightarrow[-1,1]
$$

may be written
\begin{equation}
S(\mathfrak a)
:=
\sin(q_\theta\mathfrak a).
\tag{24}
\end{equation}
In the reciprocal chart,
\begin{equation}
\boxed{
S(\mathfrak a)=\frac1\alpha.
}
\tag{25}
\end{equation}

Thus \texttt{sin()} remains ordinary \texttt{sin()}. The coordinate expression of its argument changes.

Suppress the chart map and (25) degenerates into (22).

The ``error'' is therefore informative:
\begin{equation}
\boxed{
\text{the missing projection was the hole in the notation.}
}
\tag{26}
\end{equation}

\section*{V. $\delta$ was not optional}

Equation (8) makes $\alpha$ meaningless as a complete vantage coordinate without its complement:

$$
\alpha^{-2}+\delta^{-2}=1.
$$

Neither independently determines a Euclidean decomposition without the constraint binding them.

The pair

$$
(\alpha,\delta)
$$

is therefore more fundamental to the reciprocal chart than either symbol alone.

In the circular chart:

\begin{equation}
(\alpha,\delta)
=
(\csc\theta,\sec\theta).
\tag{27}
\end{equation}

In the hyperbolic chart:

\begin{equation}
(\alpha,\delta)
=
(\coth\chi,\cosh\chi).
\tag{28}
\end{equation}

In the normalized velocity chart:

\begin{equation}
(\alpha,\delta)
=
(\beta_v^{-1},\beta_c^{-1}).
\tag{29}
\end{equation}

The same vantage relation appears respectively as circular complementarity, hyperbolic boost geometry, and reciprocal aperture.

Spatial vantage $\delta$ is therefore not an auxiliary correction to temporal vantage $\alpha$.

They are dual apertures of one constrained relation.

\section*{VI. What, then, is $\varphi$?}

If $\alpha$ and $\delta$ are reciprocal coordinates rather than ordinary circular angles, the expression

\begin{equation}
\varphi=\pi-(\alpha+\delta)
\tag{30}
\end{equation}

has no invariant meaning.

The third ``angle'' in the original triangle must therefore be retyped.

The dimensionless mixed quantity already supplied by the geometry is
\begin{equation}
m
:=
\frac{\delta}{\alpha}
=\frac{\beta_v}{\beta_c}
=\sinh\chi
=\tan\theta.
\tag{31}
\end{equation}
Hence a mixed circular angle may be recovered by

\begin{equation}
\varphi_\theta:=\arctan m=\theta,
\tag{32}
\end{equation}

while its hyperbolic coordinate is

\begin{equation}
\varphi_\chi:=\operatorname{arsinh}m=\chi.
\tag{33}
\end{equation}

Thus a third independent scalar angle is unnecessary in the one-axis case. What appeared as $\varphi$ is naturally interpreted as the \emph{transition relation} between the complementary apertures:
\begin{equation}
\boxed{
\varphi:
\alpha\leftrightarrow\delta,
\qquad
\frac{\delta}{\alpha}
=\tan\theta
=\sinh\chi.
}
\tag{34}
\end{equation}

In higher-dimensional vantage geometry, $\varphi$ may regain independent content as a mixed spatial--temporal orientation. The one-axis construction does not establish that extension; it identifies the object that such an extension must reduce to.

\section*{VII. Lift--rotate--squeeze}

A lifted Euclidean account of special relativity is otherwise vulnerable to an obvious objection:

\begin{quote}
Why lift into a larger space, rotate, and project back down if the construction was selected merely to reproduce the Lorentz transformation?
\end{quote}

The plural-angle interpretation supplies a non-circular motivation.

The physical relation admits at least the coordinate descriptions

$$
\theta,\qquad
\chi,\qquad
(\alpha,\delta).
$$

The circular representation exposes an ordinary rotation. The reciprocal representation exposes the Lorentz factors. Their relation includes an exact similarity:

$$
T_\rho=\kappa T_\beta.
$$

Hence

\begin{equation}
\boxed{
\text{lift}\rightarrow
\text{rotate}\rightarrow
\text{squeeze}
}
\tag{35}
\end{equation}

may be read as a geometric implementation of

\begin{equation}
\boxed{
\text{change vantage}\rightarrow
\text{resolve circular orientation}\rightarrow
\text{express reciprocal projection}.
}
\tag{36}
\end{equation}

The squeeze is no longer an arbitrary afterthought. A reciprocal chart demands it.

This does not by itself derive special relativity from a higher-dimensional theory. It identifies a precise mathematical structure that a successful derivation would have to explain rather than insert.

\section*{VIII. The actual claim}

No new trigonometric function is proposed.

No identity

$$
\sin x=x^{-1}
$$

is asserted for ordinary real angles.

No replacement of rapidity is required.

The proposal is only this:

\begin{equation}
\boxed{
\begin{gathered}
\text{A relativistic vantage may be treated as one geometric object}\\
\text{admitting circular, hyperbolic, and reciprocal angular charts.}
\end{gathered}}
\tag{37}
\end{equation}

For the one-axis case,
\begin{equation}
\boxed{
\begin{aligned}
\beta_v
&=\sin\theta
=\tanh\chi
=\alpha^{-1},\\
\beta_c
&=\cos\theta
=\operatorname{sech}\chi
=\delta^{-1},\\
\gamma_c
&=\csc\theta
=\coth\chi
=\alpha,\\
\gamma_v
&=\sec\theta
=\cosh\chi
=\delta,\\
\frac{\delta}{\alpha}
&=\tan\theta
=\sinh\chi.
\end{aligned}}
\tag{38}
\end{equation}

Every equality in (38) is ordinary mathematics.

The speculative step is the box around them.

\vfill
\begin{center}
\large

$$
\boxed{\texttt{sin()} \text{ stays } \texttt{sin()}.}
$$

$$
\boxed{\text{Angles become plural.}}
$$

\smallskip
\textit{Same relation. Different body.}
\end{center}

\section*{Reference}

\small
NIST Digital Library of Mathematical Functions,
\S4.23(viii), ``Gudermannian Function,'' especially
Eqs.~4.23.39--4.23.42.

\end{document}
